% GENERATED FILE. Edit canonical lessons, then run npm run generate:books.
% canonical-source-hash: fnv1a64:9020715404912777
% canonical-lessons: GU-C515-popdo, GU-C515-ukhadvu, GU-C515-kanu, GU-C515-paip, GU-C515-plambar

\chapter{Cracks and Leaks}
\label{ch:gu-cracks-and-leaks}
\hlchaptermodality{515}

\begin{chapteropening}
\textbf{By the end of this chapter:} \emph{I can describe damp, flaking walls and leaking pipes at home.}

Complete the last lesson of chapter 515: I can describe damp, flaking walls and leaking pipes at home.
\end{chapteropening}

\section[\texorpdfstring{\gu{પોપડો} (popḍo)}{પોપડો (popḍo)}]{\gu{પોપડો} (popḍo) — a flake (of plaster)}

\label{lesson:GU-C515-popdo}



\begin{quote}
Before the new one: say the Gujarati for a wall peg, then the Gujarati for to hang (something).
\end{quote}

\subsection*{You'll want to know: \gu{પોપડો}}
\textbf{\gu{પોપડો}} — \emph{popḍo} — \textquotedblleft{}a flake (of plaster)\textquotedblright{}.

\begin{quote}
\gu{દીવાલ} \gu{પરથી} \gu{એક} \gu{મોટો} \gu{પોપડો} \gu{પડ્યો}.
\end{quote}

\begin{quote}
\emph{dīvāl parthī ek moṭo popḍo paḍyo.}
\end{quote}

\begin{quote}
A big flake of plaster fell off the wall.
\end{quote}

A masculine noun.

\begin{cousinweb}
\textbf{\gu{પોપડો}} (\emph{popḍo}) comes from \emph{*pōppa}, a form CDIAL reconstructs, and CDIAL lists the Gujarati word as a detached scaly crust (CDIAL 8405).
\end{cousinweb}

\begin{grammarlens}[title={What you've built}]
One more word for describing what you lived and felt.
\end{grammarlens}

\subsection*{Your turn}
\begin{itemize}
\raggedright
  \item \emph{Say it:} \emph{popḍo}
  \item \emph{Say it:} \emph{popḍo}, once more
  \item \emph{Say it:} \emph{dīvāl parthī ek moṭo popḍo paḍyo.} — A big flake of plaster fell off the wall.
  \item \emph{Recall:} say the Gujarati for a wall peg, then the Gujarati for to hang (something), then say \emph{popḍo} again
\end{itemize}

\begin{tcolorbox}[breakable,colback=teal!4,colframe=teal!35!black,title={Before you move on}]
What does \gu{પોપડો} mean? (\textquotedblleft{}A flake (of plaster)\textquotedblright{}.) Say it once more.
\end{tcolorbox}
\section[\texorpdfstring{\gu{ઉખડવું} (ukhaḍvũ)}{ઉખડવું (ukhaḍvũ)}]{\gu{ઉખડવું} (ukhaḍvũ) — to peel off (of plaster or paint)}

\label{lesson:GU-C515-ukhadvu}



\begin{quote}
Before the new one: say the Gujarati for to hang (something), then the Gujarati for a flake (of plaster).
\end{quote}

\subsection*{You'll want to know: \gu{ઉખડવું}}
\textbf{\gu{ઉખડવું}} — \emph{ukhaḍvũ} — \textquotedblleft{}to peel off (of plaster or paint)\textquotedblright{}.

\begin{quote}
\gu{ભેજને} \gu{લીધે} \gu{છત} \gu{પરથી} \gu{પોપડા} \gu{ઉખડી} \gu{ગયા}.
\end{quote}

\begin{quote}
\emph{bhejne līdhe chat parthī popḍā ukhḍī gayā.}
\end{quote}

\begin{quote}
Because of the damp, flakes peeled off the ceiling.
\end{quote}

Past (it …): \gu{તે} \gu{પોપડો} \gu{ઉખડી} \gu{ગયો}, \emph{te popḍo ukhḍī gayo}.

\begin{cousinweb}
\textbf{\gu{ઉખડવું}} (\emph{ukhaḍvũ}) comes from \emph{*utskarati}, \textquotedblleft{}digs out\textquotedblright{}, a form CDIAL reconstructs, and CDIAL lists the Gujarati verb as \textquotedblleft{}to be rooted up\textquotedblright{} (CDIAL 1895). It is the intransitive partner of \textbf{\gu{ઉખેડવું}}, to uproot.
\end{cousinweb}

\begin{grammarlens}[title={What you've built}]
One more word for describing what you lived and felt.
\end{grammarlens}

\subsection*{Your turn}
\begin{itemize}
\raggedright
  \item \emph{Say it:} \emph{ukhaḍvũ}
  \item \emph{Say it:} \emph{ukhaḍvũ}, once more
  \item \emph{Say it:} \emph{bhejne līdhe chat parthī popḍā ukhḍī gayā.} — Because of the damp, flakes peeled off the ceiling.
  \item \emph{Recall:} say the Gujarati for to hang (something), then the Gujarati for a flake (of plaster), then say \emph{ukhaḍvũ} again
\end{itemize}

\begin{tcolorbox}[breakable,colback=teal!4,colframe=teal!35!black,title={Before you move on}]
What does \gu{ઉખડવું} mean? (\textquotedblleft{}To peel off (of plaster or paint)\textquotedblright{}.) Say it once more.
\end{tcolorbox}
\section[\texorpdfstring{\gu{કાણું} (kāṇũ)}{કાણું (kāṇũ)}]{\gu{કાણું} (kāṇũ) — a hole}

\label{lesson:GU-C515-kanu}



\begin{quote}
Before the new one: say the Gujarati for a flake (of plaster), then the Gujarati for to peel off.
\end{quote}

\subsection*{You'll want to know: \gu{કાણું}}
\textbf{\gu{કાણું}} — \emph{kāṇũ} — \textquotedblleft{}a hole\textquotedblright{}.

\begin{quote}
\gu{બાલદીમાં} \gu{એક} \gu{કાણું} \gu{છે}.
\end{quote}

\begin{quote}
\emph{bāldīmā̃ ek kāṇũ che.}
\end{quote}

\begin{quote}
There is a hole in the bucket.
\end{quote}

A neuter noun.

\begin{cousinweb}
\textbf{\gu{કાણું}} (\emph{kāṇũ}) is filed by CDIAL under Sanskrit \emph{kāṇá}, \textquotedblleft{}one-eyed\textquotedblright{}, with the Gujarati word glossed \textquotedblleft{}full of holes; a hole\textquotedblright{} (CDIAL 3019): perhaps first an empty eye-hole, as CDIAL suggests.
\end{cousinweb}

\begin{grammarlens}[title={What you've built}]
One more word for describing what you lived and felt.
\end{grammarlens}

\subsection*{Your turn}
\begin{itemize}
\raggedright
  \item \emph{Say it:} \emph{kāṇũ}
  \item \emph{Say it:} \emph{kāṇũ}, once more
  \item \emph{Say it:} \emph{bāldīmā̃ ek kāṇũ che.} — There is a hole in the bucket.
  \item \emph{Recall:} say the Gujarati for a flake (of plaster), then the Gujarati for to peel off, then say \emph{kāṇũ} again
\end{itemize}

\begin{tcolorbox}[breakable,colback=teal!4,colframe=teal!35!black,title={Before you move on}]
What does \gu{કાણું} mean? (\textquotedblleft{}A hole\textquotedblright{}.) Say it once more.
\end{tcolorbox}
\section[\texorpdfstring{\gu{પાઈપ} (pāīp)}{પાઈપ (pāīp)}]{\gu{પાઈપ} (pāīp) — a pipe}

\label{lesson:GU-C515-paip}



\begin{quote}
Before the new one: say the Gujarati for to peel off, then the Gujarati for a hole.
\end{quote}

\subsection*{You'll want to know: \gu{પાઈપ}}
\textbf{\gu{પાઈપ}} — \emph{pāīp} — \textquotedblleft{}a pipe\textquotedblright{}.

\begin{quote}
\gu{રસોડાની} \gu{પાઈપમાંથી} \gu{પાણી} \gu{ટપકે} \gu{છે}.
\end{quote}

\begin{quote}
\emph{rasoḍānī pāīpmā̃thī pāṇī ṭapke che.}
\end{quote}

\begin{quote}
Water is dripping from the kitchen pipe.
\end{quote}

A feminine noun.

\begin{cousinweb}
\textbf{\gu{પાઈપ}} (\emph{pāīp}) is borrowed from English \emph{pipe}.
\end{cousinweb}

\begin{grammarlens}[title={What you've built}]
One more word for describing what you lived and felt.
\end{grammarlens}

\subsection*{Your turn}
\begin{itemize}
\raggedright
  \item \emph{Say it:} \emph{pāīp}
  \item \emph{Say it:} \emph{pāīp}, once more
  \item \emph{Say it:} \emph{rasoḍānī pāīpmā̃thī pāṇī ṭapke che.} — Water is dripping from the kitchen pipe.
  \item \emph{Recall:} say the Gujarati for to peel off, then the Gujarati for a hole, then say \emph{pāīp} again
\end{itemize}

\begin{tcolorbox}[breakable,colback=teal!4,colframe=teal!35!black,title={Before you move on}]
What does \gu{પાઈપ} mean? (\textquotedblleft{}A pipe\textquotedblright{}.) Say it once more.
\end{tcolorbox}
\section[\texorpdfstring{\gu{પ્લમ્બર} (plambar)}{પ્લમ્બર (plambar)}]{\gu{પ્લમ્બર} (plambar) — a plumber}

\label{lesson:GU-C515-plambar}



\begin{quote}
Before the new one: say the Gujarati for a hole, then the Gujarati for a pipe.
\end{quote}

\subsection*{You'll want to know: \gu{પ્લમ્બર}}
\textbf{\gu{પ્લમ્બર}} — \emph{plambar} — \textquotedblleft{}a plumber\textquotedblright{}.

\begin{quote}
\gu{પ્લમ્બરે} \gu{રસોડાની} \gu{જૂની} \gu{પાઈપ} \gu{બદલી}.
\end{quote}

\begin{quote}
\emph{plambare rasoḍānī jūnī pāīp badlī.}
\end{quote}

\begin{quote}
The plumber changed the old kitchen pipe.
\end{quote}

A masculine or feminine noun, by the person meant.

\begin{cousinweb}
\textbf{\gu{પ્લમ્બર}} (\emph{plambar}) is borrowed from English \emph{plumber}.
\end{cousinweb}

\begin{grammarlens}[title={What you've built}]
One more word for describing what you lived and felt.
\end{grammarlens}

\subsection*{Your turn}
\begin{itemize}
\raggedright
  \item \emph{Say it:} \emph{plambar}
  \item \emph{Say it:} \emph{plambar}, once more
  \item \emph{Say it:} \emph{plambare rasoḍānī jūnī pāīp badlī.} — The plumber changed the old kitchen pipe.
  \item \emph{Recall:} say the Gujarati for a hole, then the Gujarati for a pipe, then say \emph{plambar} again
\end{itemize}

\begin{tcolorbox}[breakable,colback=teal!4,colframe=teal!35!black,title={Before you move on}]
What does \gu{પ્લમ્બર} mean? (\textquotedblleft{}A plumber\textquotedblright{}.) Say it once more.
\end{tcolorbox}
